% ============================================================================
% Fichier  : partie1/P01_C03_gouvernance.tex
% Titre    : Gouvernance des Risques
% Auteur   : MINKA MI NGUIDJOÏ Thierry Emmanuel
% Version  : 1.0 -- novembre 2026
% ============================================================================
% Sources : IIA, modèle des trois lignes (2020) ; COSO ERM (2017) ;
% ISO 31000:2018 (leadership et engagement). Organisation de LiZbiZ fictive.
% ============================================================================

\chapter{Gouvernance des Risques}
\label{chap:gouvernance}
\chaptermark{Gouvernance des risques}

\epigraph{%
    « Un risque sans propriétaire\\
    est un risque que personne ne traite. »%
}{Principe du cours}

\niveauChapitre{Fondamental}{Chapitres~\ref{chap:risque}
et~\ref{chap:perception}.}

\begin{boxobjectifs}
\begin{itemize}
    \item Répartir les responsabilités entre conseil, direction, propriétaires
          de risques, fonction risques et audit interne.
    \item Appliquer le modèle des trois lignes de l'IIA.
    \item Formuler une appétence et des tolérances au risque.
    \item Définir le rôle et la composition d'un comité des risques.
    \item Maîtriser les risques d'une activité externalisée.
\end{itemize}
\end{boxobjectifs}

\bigskip

% ============================================================================
\section*{Question directrice : qui, chez LiZbiZ, répondait du risque du prestataire~?}
% ============================================================================

Le contrat avec SAHEL INFOGÉRANCE avait été signé par la direction des
opérations, négocié par la direction financière, et suivi par
l'informatique. Chacun pensait que l'un des deux autres veillait à la
sécurité. Personne n'en était chargé par écrit. Ce chapitre montre comment
attribuer chaque risque à quelqu'un, et vérifier qu'il s'en occupe.

\medskip
\textbf{Plan du chapitre.}
\begin{enumerate}
    \item Les rôles (section~\ref{sec:c03-roles}).
    \item Le modèle des trois lignes (section~\ref{sec:c03-lignes}).
    \item Appétence et tolérance (section~\ref{sec:c03-appetence}).
    \item Le comité des risques (section~\ref{sec:c03-comite}).
    \item Les risques externalisés (section~\ref{sec:c03-tiers}).
\end{enumerate}

\separateur

% ============================================================================
\section{Les Rôles}
\label{sec:c03-roles}
% ============================================================================

\begin{boxdef}
\termecle{Gouvernance des risques}~: ensemble des rôles, des règles et des
instances par lesquels une organisation fixe le niveau de risque qu'elle
accepte, attribue chaque risque à un responsable et vérifie que les
risques sont traités.
\end{boxdef}

\begin{tableaucours}{L{3.4cm}L{8.8cm}}{%
    \tableauHeader{Acteur} &
    \tableauHeader{Rôle}
}
Conseil d'administration & Fixe l'appétence au risque, approuve la
politique, reçoit le compte rendu des risques majeurs \\
Direction générale & Met en œuvre la politique, arbitre les ressources,
répond devant le conseil \\
Propriétaire de risque & Responsable de l'activité exposée~; décide du
traitement de son risque et en rend compte \\
Fonction risques & Fournit la méthode, consolide le registre, alerte,
conseille~; ne possède pas les risques \\
Audit interne & Donne au conseil une assurance indépendante sur le
dispositif \\
\end{tableaucours}
\vspace{-0.8em}
\begin{center}{\small\itshape Tableau~3.1. Les rôles de la gouvernance des risques}\end{center}

\begin{boxalert}
\textbf{Violation typique.} Faire de la fonction risques (ou du service
informatique) le propriétaire de tous les risques.

\smallskip
\textbf{Règle~:} le propriétaire d'un risque est celui qui a l'autorité et
le budget pour le traiter, c'est-à-dire le responsable de l'activité
exposée. La fonction risques l'aide, elle ne le remplace pas.
\end{boxalert}

% ============================================================================
\section{Le Modèle des Trois Lignes}
\label{sec:c03-lignes}
% ============================================================================

L'Institut des auditeurs internes (IIA) a publié en 2020 une mise à jour
du modèle des « trois lignes »~\cite{iia2020troislignes}. Il décrit les rôles
plutôt que des barrières~: la première ligne gère les risques dans les
opérations, la deuxième apporte expertise, méthode et surveillance, la
troisième, l'audit interne, donne une assurance indépendante à l'organe de
gouvernance.

\begin{center}
\begin{tikzpicture}[
    b/.style = {draw=CouleurBordDef, fill=CouleurDef, rounded corners=3pt,
        minimum width=3.4cm, minimum height=1.0cm, align=center,
        font=\small\color{CouleurPrimaire}},
    arrow/.style = {-{Stealth[length=4pt]}, thick, color=CouleurSecondaire}]
\node[b, fill=CouleurPrimaire, text=white, minimum width=11cm] (g) at (3.9, 2.2)
    {\color{white}\bfseries Organe de gouvernance (conseil)};
\node[b] (l1) at (0, 0) {\textbf{1\textsuperscript{re} ligne}\\opérations, métiers};
\node[b] (l2) at (3.9, 0) {\textbf{2\textsuperscript{e} ligne}\\risques, conformité};
\node[b] (l3) at (7.8, 0) {\textbf{3\textsuperscript{e} ligne}\\audit interne};
\draw[arrow] (l1) -- (g.south -| l1);
\draw[arrow] (l2) -- (g.south -| l2);
\draw[arrow] (l3) -- (g.south -| l3);
\draw[arrow, dashed] (l2) -- (l1);
\end{tikzpicture}
\end{center}
\vspace{-0.4em}
\begin{center}{\small\itshape Figure~3.1. Le modèle des trois lignes (d'après l'IIA, 2020)}\end{center}

\begin{boxscenario}{Opération LIZBIZ. La nouvelle organisation}
\textbf{Constatation} (avril 2026). LiZbiZ recrute une gestionnaire des
risques, Mme~MOUAFO, rattachée au directeur général (deuxième ligne). Mme
NKENG, contrôleuse interne, voit son rôle d'audit confirmé et rattaché au
comité d'audit du conseil (troisième ligne). Chaque directeur devient
propriétaire des risques de son activité (première ligne).

\medskip
\textbf{Analyse} (confiance élevée). La séparation entre la fonction qui
conseille (Mme~MOUAFO) et celle qui contrôle (Mme~NKENG) garantit que
l'audit ne juge pas un dispositif qu'il a lui-même conçu.
\end{boxscenario}

% ============================================================================
\section{Appétence et Tolérance}
\label{sec:c03-appetence}
% ============================================================================

\begin{boxdef}
\termecle{Appétence au risque} (COSO ERM)~: types et niveau de risque qu'une
organisation est disposée à accepter pour atteindre sa
mission~\cite{coso2017erm}. La \termecle{tolérance} précise, pour un objectif
donné, l'écart acceptable autour de la cible, sous forme mesurable.
\end{boxdef}

\begin{tableaucours}{L{4.0cm}L{8.2cm}}{%
    \tableauHeader{Domaine} &
    \tableauHeader{Énoncé adopté par le conseil de LiZbiZ (mai 2026)}
}
Continuité & Aucun arrêt de la facturation de plus de deux jours ouvrés \\
Données clients & Aucune fuite de données personnelles acceptée~; tout
traitement de ces données est chiffré \\
Pertes financières & Perte annuelle liée aux incidents informatiques
inférieure à 100~millions de FCFA dans neuf ans sur dix \\
Prestataires & Aucun accès distant d'un tiers sans second facteur ni
journalisation \\
\end{tableaucours}
\vspace{-0.8em}
\begin{center}{\small\itshape Tableau~3.2. L'appétence au risque de LiZbiZ}\end{center}

Une bonne appétence est \emph{mesurable}~: on doit pouvoir dire, à la fin
de l'année, si elle a été respectée. « Nous avons une faible appétence au
risque cyber » ne se vérifie pas~; « aucun arrêt de plus de deux jours »
se vérifie.

% ============================================================================
\section{Le Comité des Risques}
\label{sec:c03-comite}
% ============================================================================

\begin{boxprocedure}
\textbf{Le comité des risques de LiZbiZ}
\begin{enumerate}
    \item \textbf{Composition}~: directeur général (président), directeurs
          propriétaires de risques, gestionnaire des risques (secrétaire),
          contrôleuse interne (invitée, sans voix délibérative).
    \item \textbf{Fréquence}~: trimestrielle, et extraordinaire après tout
          incident majeur.
    \item \textbf{Ordre du jour}~: évolution des risques majeurs, état des
          plans de traitement, indicateurs, dérogations, risques
          émergents.
    \item \textbf{Décisions}~: consignées avec le propriétaire, l'échéance
          et le budget.
    \item \textbf{Compte rendu} au conseil une fois par semestre.
\end{enumerate}
\end{boxprocedure}

% ============================================================================
\section{Les Risques Externalisés}
\label{sec:c03-tiers}
% ============================================================================

Externaliser une activité ne transfère pas la responsabilité de ses
risques. LiZbiZ reste exploitant de ses systèmes au sens de la loi
(Chapitre~\ref{chap:droit}), même si SAHEL INFOGÉRANCE les administre.

\begin{tableaucours}{L{3.4cm}L{8.8cm}}{%
    \tableauHeader{Étape} &
    \tableauHeader{Exigence}
}
Avant le contrat & Évaluer le prestataire~; préciser ses accès et ses
obligations de sécurité \\
Contrat & Délai de notification d'un incident, second facteur, droit
d'audit, réversibilité, assurance \\
Pendant & Revue périodique des accès, indicateurs, audit \\
Fin & Retrait des accès, restitution et effacement des données \\
\end{tableaucours}
\vspace{-0.8em}
\begin{center}{\small\itshape Tableau~3.3. Maîtriser un prestataire}\end{center}

\begin{boxobjet}{R-02 · Accès distant d'un prestataire}
\textbf{Scénario~:} un compte du prestataire est usurpé et donne accès au
domaine. \textbf{Cotation~:} vraisemblance 3, impact 4 (score~12,
intolérable). \textbf{Traitement~:} réduire (second facteur, bastion,
plages horaires) et contractualiser (notification sous 24~heures, droit
d'audit). \textbf{Propriétaire~:} M.~NGUELE, signataire du contrat.
\textbf{Indicateur~:} part des accès du prestataire passés par le bastion.
\end{boxobjet}

\begin{boxmetier}{Auditeur des systèmes d'information}
\textbf{Vérifier que chaque risque a un nom.} L'auditeur prend le registre
et, pour chaque risque majeur, vérifie que le propriétaire désigné sait
qu'il l'est, connaît son plan de traitement et peut montrer où il en est.
\end{boxmetier}

\separateur

\begin{boxresume}
\textbf{Ce qu'il faut retenir de ce chapitre~:}
\begin{itemize}
    \item \textbf{Le conseil} fixe l'appétence~; \textbf{la direction} met
          en œuvre~; \textbf{le propriétaire} traite son risque.
    \item \textbf{Trois lignes}~: les opérations gèrent, la fonction risques
          aide et surveille, l'audit donne une assurance indépendante.
    \item \textbf{L'appétence} se formule de manière mesurable.
    \item \textbf{Le comité des risques} décide, consigne, rend compte.
    \item \textbf{Externaliser} ne transfère pas la responsabilité.
\end{itemize}
\end{boxresume}

\begin{boxinfo}
\textbf{Points de contrôle de l'élève.} Je sais~: attribuer un risque à un
propriétaire~; placer un acteur dans les trois lignes~; rédiger un énoncé
d'appétence mesurable~; décrire un comité des risques~; énumérer les
clauses de sécurité d'un contrat d'externalisation.
\end{boxinfo}

\bigskip

\begin{boxexo}[title={\ding{45}\quad Exercice~3.1. Quelle ligne~? \niveaubase}]
Placez dans les trois lignes~: le chef d'agence de Garoua~; Mme~MOUAFO~;
Mme~NKENG~; le RSSI~; un auditeur externe mandaté par le conseil.
\end{boxexo}

\begin{boxexo}[title={\ding{45}\quad Exercice~3.2. Le propriétaire \niveaubase}]
Qui doit être propriétaire du risque « fraude au virement »~: la DAF, le
RSSI ou la gestionnaire des risques~? Justifiez.
\end{boxexo}

\begin{boxexo}[title={\ding{45}\quad Exercice~3.3. Mesurable \niveaubase}]
Réécrivez de manière mesurable~: « LiZbiZ a une faible tolérance aux
interruptions de service ».
\end{boxexo}

\begin{boxexo}[title={\ding{45}\quad Exercice~3.4. Indépendance \niveaumoyen}]
Pourquoi Mme~NKENG ne doit-elle pas rédiger elle-même la méthode d'analyse
des risques qu'elle auditera~?
\end{boxexo}

\begin{boxexo}[title={\ding{45}\quad Exercice~3.5. Le contrat \niveaumoyen}]
Rédigez cinq clauses de sécurité à ajouter au contrat de SAHEL
INFOGÉRANCE.
\end{boxexo}

\begin{boxexo}[title={\ding{45}\quad Exercice~3.6. Appétence et décision \niveaumoyen}]
Une option coûte 30~millions et ramène la perte dans neuf ans sur dix à
80~millions~; une autre coûte 10~millions et la ramène à 120~millions.
Laquelle respecte l'appétence du Tableau~3.2~?
\end{boxexo}

\begin{boxexo}[title={\ding{45}\quad Exercice~3.7. L'ordre du jour \niveaumoyen}]
Rédigez l'ordre du jour du comité des risques extraordinaire du 12~mars.
\end{boxexo}

\begin{boxexo}[title={\ding{45}\quad Exercice~3.8. RACI \niveaumoyen}]
Établissez une matrice RACI (réalise, approuve, consulté, informé) pour la
mise en place du second facteur sur la passerelle.
\end{boxexo}

\begin{boxexo}[title={\ding{45}\quad Exercice~3.9. Charte \niveauexpert}]
Rédigez la charte de la fonction risques de LiZbiZ~: mission, rattachement,
droits d'accès à l'information, relations avec l'audit.
\end{boxexo}

\begin{boxexo}[title={\ding{45}\quad Exercice~3.10. Une banque \niveauexpert}]
Comment la gouvernance des risques de la BANQUE SANAGA doit-elle différer
de celle de LiZbiZ~? Considérez la taille, la supervision et la nature
des risques.
\end{boxexo}

\bigskip

\begin{boxinfo}
\textbf{Pour aller plus loin.}
\begin{itemize}
    \item IIA, \emph{The IIA's Three Lines Model}~\cite{iia2020troislignes}.
    \item COSO, \emph{Enterprise Risk Management}~\cite{coso2017erm}.
\end{itemize}
\end{boxinfo}

\begin{boxinfo}
\textbf{Transition.} Les rôles posés, il faut une méthode commune. Le
Chapitre~\ref{chap:cadres} présente l'ISO~31000 et le COSO ERM.
\end{boxinfo}
